\section*{M5: Why growing tree size is out of scope without further hypotheses}

The theorem in \texttt{fixed\_tree\_convergence.tex} fixes $T$ (hence
$k$) independent of $N$. This section records, precisely, why dropping
that restriction is unsound without additional summability hypotheses —
so that no future session is tempted to silently generalize the theorem
to $k=k_N\to\infty$.

\textbf{The telescoping sum has $k_N$ terms at level $N$.} Revisiting the
inductive step: the error accumulated at the root is (informally) a sum
over all $k_N$ internal vertices of a per-vertex error term
$\varepsilon_{v,N}$ (from H1/H2), each individually $\to 0$ as $N\to
\infty$ \emph{for fixed $v$}. If $k_N\to\infty$, the total error is a sum
of a \emph{growing number} of terms, each shrinking — whether the total
shrinks, stays bounded, or diverges depends entirely on the
\emph{relative rates}: a sum of $k_N$ terms each of size
$O(k_N^{-1-\delta})$ (for any $\delta>0$) still $\to 0$, but a sum of
$k_N$ terms each of size $\Theta(k_N^{-1})$ does \emph{not} — it stays
$\Theta(1)$, contradicting convergence to $0$. No hypothesis in
(H1)–(H4) constrains this rate; they only assert convergence at each
fixed $v$, individually.

\textbf{Minimal illustrative counterexample (schematic, not a
constructed instance in this repository).} Take a sequence of trees
$T_N$ with $k_N=N$ internal vertices (a growing chain), and let each
node's discretization error satisfy $\varepsilon_{v,N} = c/N$ for a fixed
constant $c>0$ — this is consistent with (H1) holding at every fixed
node (as $N\to\infty$ with $v$ fixed, $\varepsilon_{v,N}\to0$), but the
accumulated root error is $\sum_{v=1}^{N} c/N = c$, a nonzero constant,
for every $N$ — it does \emph{not} converge to $0$. This shows (H1)–(H4)
alone are insufficient once $k$ grows with $N$; a genuine result for
growing trees needs an explicit summability condition (e.g.\
$\sum_N \sup_v \varepsilon_{v,N} < \infty$, or a depth/rate tradeoff),
which this pass does not attempt to state or prove.

\textbf{Status.} Not a counterexample to the fixed-tree theorem (which
makes no claim about growing $k$) — a documented boundary of its scope,
per the mission's own explicit instruction not to claim results for
growing tree size without controlling depth, number of sources, and
summability.
